\documentclass[11pt,a4paper]{article}
\usepackage[utf8]{inputenc}
\usepackage[T1]{fontenc}
\usepackage{amsmath,amssymb}
\usepackage{booktabs}
\usepackage{geometry}
\usepackage{enumitem}
\usepackage{hyperref}

\geometry{margin=2.5cm}

\title{Types Won't Save You, But Typestates Will: Static Protocol\\
Checking for C++ gRPC}
\author{Mihai Tira}
\date{October 2026}

\begin{document}

\maketitle

\section{Research Report Outline}

This report studies static verification of synchronous C++ gRPC streams against Binary Session Type contracts. It details the end-to-end verification pipeline: extracting binary session types from Protobuf custom options, compiling them into dual finite-state automata via our session compiler, and verifying synchronous C++ gRPC streams using a path-sensitive Clang AST static verifier. Native in-tree compiler-plugin integration remains future work.

\begin{enumerate}
    \item \textbf{Introduction}
    \begin{itemize}
        \item Motivation: protocol safety in distributed streaming RPCs
        \item Problem: message typing does not capture temporal or payload constraints
        \item Research questions, hypothesis, scope, and contributions
    \end{itemize}

    \item \textbf{Background and Related Work}
    \begin{itemize}
        \item Binary and multiparty session types, duality, and subtyping~\cite{honda-v-k-1998,yoshida-vasconcelos-2007,honda-yoshida-carbone-2008,honda-et-al-2011,gay-hole-2005}
        \item Typestate, type-and-effect systems, and refinement session types~\cite{strom-yemini-1986,nielson-nielson-1999,orchard-yoshida-2016,kouzapas-et-al-2016,pradel-et-al-2012,ueno-das-2026}
        \item Protocol verification in C, C++, and other languages~\cite{hu-yoshida-2016,ng-et-al-2012,ng-yoshida-honda-2012-cc,marques-et-al-2013,alsubhi-et-al-2025}
        \item Comparison with existing approaches and research gap
    \end{itemize}

    \item \textbf{Formal Model}
    \begin{itemize}
        \item Session-type grammar and endpoint duality
        \item Subtyping and stream termination semantics
        \item Two levels of contracts: call sequence and lifecycle; payload variants and field-presence conditions
        \item Translation of session types into deterministic finite-state machines
    \end{itemize}

    \item \textbf{Static Verification}
    \begin{itemize}
        \item Verification pipeline and Clang AST name resolution
        \item Path-sensitive control-flow analysis, including read results and early returns
        \item Payload tracking and bounded loop analysis
    \end{itemize}

    \item \textbf{Empirical Evaluation}
    \begin{itemize}
        \item Repository audit and analyzer evaluation
        \item Real-world cases from WeNet, Apache Arrow Flight, BuildBox, gRPC, Triton, and Google Cloud C++
        \item Mutation testing across seven violation categories
        \item Annotation effort, analysis time, and threats to validity
    \end{itemize}

    \item \textbf{Conclusion, Limitations and Future Work}
    \begin{itemize}
        \item Answers to the research questions
        \item Automated specification generation and native compiler integration
    \end{itemize}
\end{enumerate}

\section{Applied Research Plan}

\subsection{Working Hypothesis}
Call ordering, stream termination, and Protobuf payload constraints in synchronous C++ gRPC can be represented as dual session-type automata and checked through path-sensitive Clang AST analysis, without changing application code or generated headers.

\subsection{Methodology}
\begin{enumerate}
    \item \textbf{Repository audit:} Identify lifecycle and payload rules across open-source C++ gRPC codebases.
    \item \textbf{Contract formalization:} Define dual client and server state machines, including payload conditions.
    \item \textbf{Verifier implementation:} Build path-sensitive analysis for API binding, early returns, and bounded loop unrolling.
    \item \textbf{Evaluation:} Test real-world cases and a multi-category mutation benchmark, measuring precision, recall, and false positives.
\end{enumerate}

\subsection{Original Approach}
Unlike approaches that require custom communication libraries, generated state wrappers, or invasive template changes, this method checks standard C++ gRPC code directly. External contracts describe expected behavior, while Clang AST analysis tracks stream states and payload changes across branches, early returns, and loops.

\subsection{Planned Experiments}
\begin{enumerate}
    \item \textbf{Real-world verification:} Analyze multiple cases from the repositories and compare findings with the repositories' higher-level abstractions.
    \item \textbf{Fault injection:} Measure precision and recall for valid behavior, incorrect initialization, missing write completion, early-return leaks, writes after closure, invalid payloads, and premature termination.
    \item \textbf{Overhead:} Measure the need for source annotations and analysis time per translation unit.
\end{enumerate}

\section{Original Contributions and Research Questions}

\subsection{Comparison with Existing Work}
Table~\ref{tab:comparison} summarizes the main differences between this approach and related systems.

\begin{table}[htbp]
\centering
\small
\begin{tabular}{@{}p{2.3cm}p{1.6cm}p{2.5cm}p{2.0cm}p{2.3cm}p{2.5cm}@{}}
\toprule
\textbf{System} & \textbf{Domain} & \textbf{Specification} & \textbf{API changes} & \textbf{Payload tracking} & \textbf{Evaluation} \\
\midrule
Session C~\cite{ng-yoshida-honda-2012-cc} & C / MPI & External Scribble & Yes & Branch labels & Rewritten HPC kernels \\
Marques et al.~\cite{marques-et-al-2013} & C / MPI & VCC contracts & No; annotations & Rank and length & Small MPI examples \\
Mungo~\cite{kouzapas-et-al-2016} & Java & Scribble and typestate & Generates skeletons & Enum returns & SMTP client \\
Hu and Yoshida~\cite{hu-yoshida-2016} & Java & Generated state classes & Yes & Method per label & Microbenchmarks \\
Coconut~\cite{alsubhi-et-al-2025} & C++ & Template DSL & Yes & Primitive returns & Ported drivers and mutations \\
\textbf{This study} & \textbf{C++ gRPC} & \textbf{External dual state machines} & \textbf{No} & \textbf{Payload variants and field presence} & \textbf{Unmodified open-source code} \\
\bottomrule
\end{tabular}
\caption{Comparison with related session-type and typestate verification approaches.}
\label{tab:comparison}
\end{table}

\subsection{Original Contributions}
\begin{enumerate}
    \item \textbf{End-to-end IDL-to-compiler pipeline:} A binary session type compiler extracts protocol options from \texttt{.proto} definitions, computes endpoint duality ($\overline{S}$), and synthesizes dual finite-state automata (\texttt{grpc.session.v1}).
    \item \textbf{Verification without code changes:} A path-sensitive Clang AST analyzer checks standard C++ gRPC stream APIs against the compiled automata.
    \item \textbf{Payload-aware analysis:} The verifier tracks Protobuf variants, field clearing, early returns, and read-result branches, with bounded loop analysis.
    \item \textbf{Real-world evaluation:} The approach is assessed on existing open-source repositories, alongside an audit of their higher-level abstractions.
\end{enumerate}

\subsection{Research Questions}
\begin{itemize}
    \item \textbf{RQ1 --- Expressiveness:} How well can dual session-type automata capture temporal and payload constraints in synchronous C++ gRPC streams?
    \item \textbf{RQ2 --- Defect detection:} Can path-sensitive AST analysis find protocol and lifecycle defects in existing C++ gRPC repositories?
    \item \textbf{RQ3 --- Precision:} Can it distinguish genuine violations from safe patterns such as RAII wrappers, transport-failure exits, and loop-carried field clearing?
    \item \textbf{RQ4 --- Practicality:} What annotation effort and analysis time does the approach require compared with existing C/C++ protocol checkers?
\end{itemize}

\bibliographystyle{IEEEtran}
\bibliography{references}

\end{document}
